% ============================================================================
% Fichier  : partie3/P03_C08_constater_rapporter_suivre.tex
% Titre    : Mesurer, Conclure, Rapporter, Suivre
% Auteur   : MINKA MI NGUIDJOÏ Thierry Emmanuel
% Version  : 0.2 -- novembre 2026 (réécriture étoffée, alignée sur INF3226
%            C11 et C12, A0 et A1)
% ============================================================================
% Sources : A1 (hal-05701270) : mu, U, H, cas limites, quadrant, divergence ;
% A0 (hal-05701007) : non-conformité connue / inconnue, axiomes 5 et 6 ;
% Shannon (1948), Kullback et Leibler (1951), Cover et Thomas (2006).
% Faits MINÉQUIR (répondants, mars 2026, 20 réponses par domaine) :
%   D1 Accès 11/3/6 (0,55 ; 0,30 ; 1,41) conformité opaque ;
%   D2 Fournisseurs 6/12/2 (0,30 ; 0,10 ; 1,30) non-conformité claire ;
%   D3 Circuit du mandat 15/3/2 (0,75 ; 0,10 ; 1,05) conformité claire ;
%   D4 Exploitation 8/4/8 (0,40 ; 0,40 ; 1,52) non-conformité opaque.
% Audit de suivi de novembre 2026, D1 : 17/2/1 (0,85 ; 0,05 ; 0,75) ;
%   D_KL = 0,346 bit (0,247 avec lissage de Laplace).
% Rapport provisoire 17/04/2026 ; réponse DAF 30/04/2026 ; définitif
% 15/05/2026 ; recommandations R-01 à R-06 (LAB-10).
% ============================================================================

\chapter[Mesurer, conclure, rapporter, suivre]{Mesurer, Conclure, Rapporter, Suivre}
\label{chap:constater}
\chaptermark{Conclure, rapporter, suivre}

\epigraph{%
    « On ne lit jamais l'entropie sans l'incertitude,\\
    ni l'incertitude sans la conformité. »%
}{Règle du cours, d'après A1}

\niveauChapitre{Intermédiaire}{Chapitre~\ref{chap:preuve}.}

\begin{boxobjectifs}
\begin{itemize}
    \item Calculer et interpréter les trois mesures $\mu$, $U$ et $H$ d'un
          domaine, et en connaître les cas limites.
    \item Décomposer l'entropie en part d'ignorance et part de désaccord.
    \item Situer un domaine dans le quadrant conformité~$\times$~incertitude
          et en déduire l'action.
    \item Appliquer la règle de décision du cours~: $U$ d'abord, $\mu$
          ensuite, $H$ pour contrôler.
    \item Rédiger une observation en sept éléments, et la relier au constat
          en cinq éléments.
    \item Conduire la procédure contradictoire et rédiger le rapport.
    \item Suivre les recommandations et comparer deux cycles par la
          divergence de Kullback-Leibler.
    \item Préserver la preuve quand le constat peut fonder une suite.
\end{itemize}
\end{boxobjectifs}

\bigskip

% ============================================================================
\section*{Pourquoi ce chapitre}
% ============================================================================

Les cotes \OUI{}, \NON{} et \NEK{} sont des résultats de travail. Pour
conclure, il faut les lire ensemble, domaine par domaine, sans confondre ce
que l'on sait et ce que l'on ignore. L'article A1 donne pour cela trois
mesures et un quadrant~\cite{minka2026nek}~; l'article A0, la distinction
entre non-conformité connue et non-conformité
inconnue~\cite{minka2026auditologie}. Ce chapitre les applique au dossier
MINÉQUIR, puis suit le rapport jusqu'au plan d'action et au cycle suivant.

\medskip
\textbf{Plan du chapitre.}
\begin{enumerate}
    \item Trois mesures (section~\ref{sec:c08-mesures}) et la décomposition de
          l'entropie (section~\ref{sec:c08-decomp}).
    \item Le quadrant et la règle de décision (sections~\ref{sec:c08-quadrant}
          et~\ref{sec:c08-regle}).
    \item Le constat et l'observation (section~\ref{sec:c08-constat}).
    \item Le rapport et la procédure contradictoire
          (sections~\ref{sec:c08-rapport} et~\ref{sec:c08-contradictoire}).
    \item Suivre les recommandations et comparer les cycles
          (sections~\ref{sec:c08-suivi} et~\ref{sec:c08-cycles}).
    \item Quand le constat devient une affaire (section~\ref{sec:c08-suites}).
\end{enumerate}

\separateur

% ============================================================================
\section{Trois mesures}
\label{sec:c08-mesures}
% ============================================================================

\subsection{Les définitions}

Pour un ensemble de réponses (une question posée à plusieurs répondants, ou
toutes les questions d'un domaine), on note $p(\OUI)$, $p(\NON)$ et
$p(\NEK)$ les parts de chaque réponse.

\begin{boxdef}
\termecle{Taux de conformité}\index{taux de conformité}~: $\mu = p(\OUI)$.

\medskip
\termecle{Taux d'incertitude}\index{taux d'incertitude}~: $U = p(\NEK)$.

\medskip
\termecle{Entropie}\index{entropie}~: dispersion des réponses, en bits,
\[
    H = -\sum_{m \in \{\OUI, \NON, \NEK\}} p(m) \log_2 p(m),
    \qquad 0 \leq H \leq \log_2 3 \approx 1{,}585,
\]
avec la convention $0 \log_2 0 = 0$.
\end{boxdef}
{\small\itshape Définitions de l'article A1~\cite{minka2026nek}~;
l'entropie est celle de Shannon~\cite{shannon1948}.}

\medskip

L'entropie est nulle quand tous répondent la même chose, maximale quand les
trois réponses sont également fréquentes. Au niveau d'un domaine, on agrège
les réponses de toutes ses questions avant de calculer~: les mesures se
recalculent toujours à partir des réponses, jamais comme moyenne des
mesures du niveau inférieur.

\subsection{Quatre cas limites}

\begin{tableaucours}{L{4.6cm}cccL{3.4cm}}{%
    \tableauHeader{Situation} &
    \tableauHeader{$\mu$} & \tableauHeader{$U$} & \tableauHeader{$H$} &
    \tableauHeader{Lecture}
}
Tous répondent \OUI{} & 1 & 0 & 0 & Conformité certaine \\
Tous répondent \NON{} & 0 & 0 & 0 & Non-conformité certaine \\
Tous répondent \NEK{} & 0 & 1 & 0 & \textbf{Ignorance unanime} \\
Réponses également réparties & $1/3$ & $1/3$ & 1,585 & Dispersion maximale \\
\end{tableaucours}
\vspace{-0.8em}
\begin{center}{\small\itshape Tableau~8.1. Les cas limites des trois
mesures (d'après~\cite{minka2026nek})}\end{center}

\begin{boxalert}
\textbf{Violation typique.} Conclure «~certitude~» parce que l'entropie est
nulle.

\smallskip
\textbf{Règle~:} une entropie nulle signifie seulement que tous ont donné la
même réponse. Si cette réponse est \NEK{}, personne ne sait rien~: $H = 0$ et
$U = 1$. On ne lit jamais $H$ sans $U$ ni $\mu$.
\end{boxalert}

\subsection{Les quatre domaines du MINÉQUIR}

En mars~2026, la vérificatrice a posé cinq questions par domaine aux quatre
interlocuteurs du Chapitre~\ref{chap:preuve}, soit vingt réponses par
domaine.

\begin{tableaucours}{L{3.6cm}cccccc}{%
    \tableauHeader{Domaine} &
    \tableauHeader{\OUI{}} & \tableauHeader{\NON{}} & \tableauHeader{\NEK{}} &
    \tableauHeader{$\mu$} & \tableauHeader{$U$} & \tableauHeader{$H$}
}
D1 Accès & 11 & 3 & 6 & 0,55 & 0,30 & 1,41 \\
D2 Fichier des fournisseurs & 6 & 12 & 2 & 0,30 & 0,10 & 1,30 \\
D3 Circuit du mandat & 15 & 3 & 2 & 0,75 & 0,10 & 1,05 \\
D4 Exploitation & 8 & 4 & 8 & 0,40 & 0,40 & 1,52 \\
\end{tableaucours}
\vspace{-0.8em}
\begin{center}{\small\itshape Tableau~8.2. Les trois mesures des quatre
domaines (réponses des répondants, mars~2026)}\end{center}

\begin{boxoutils}{Calculer les trois mesures}
Dans un tableur, avec les effectifs en B2 (\OUI{}), C2 (\NON{}), D2
(\NEK{})~:
\begin{lstlisting}
E2 (mu) : =B2/SOMME(B2:D2)
F2 (U)  : =D2/SOMME(B2:D2)
G2 (H)  : =-SOMME(SI(B2:D2>0;B2:D2/SOMME(B2:D2)*LOG(B2:D2/SOMME(B2:D2);2);0))
\end{lstlisting}
La dernière formule se valide comme formule matricielle dans les tableurs
qui l'exigent.
\end{boxoutils}

\separateur

% ============================================================================
\section{Décomposer l'entropie}
\label{sec:c08-decomp}
% ============================================================================

Une entropie élevée peut venir de deux choses très différentes~: beaucoup de
répondants ne savent pas (ignorance), ou ceux qui savent ne sont pas
d'accord (désaccord entre \OUI{} et \NON{}). On ne traite pas l'ignorance
comme le désaccord.

On regroupe les réponses en deux classes~: \NEK{} (probabilité $U$) et
«~réponse informée~» (\OUI{} ou \NON{}, probabilité $1 - U$). Dans la seconde
classe, la part des \OUI{} vaut $\mu / (1 - U)$. La règle de chaînage de
l'entropie~\cite{shannon1948,cover2006elements} donne, pour $U < 1$~:
\[
    H = \underbrace{H_b(U)}_{\text{ignorance}} +
        \underbrace{(1 - U)\, H_b\!\left(\frac{\mu}{1 - U}\right)}_{\text{désaccord}},
    \qquad H_b(x) = -x\log_2 x - (1-x)\log_2(1-x).
\]

\begin{tableaucours}{L{3.6cm}ccccL{2.8cm}}{%
    \tableauHeader{Domaine} &
    \tableauHeader{$H$} &
    \tableauHeader{$H_b(U)$} &
    \tableauHeader{Désaccord} &
    \tableauHeader{Part d'ignorance} &
    \tableauHeader{Lecture}
}
D1 Accès & 1,41 & 0,88 & 0,53 & 63~\% & Surtout ignorance \\
D2 Fournisseurs & 1,30 & 0,47 & 0,83 & 36~\% & Surtout désaccord \\
D4 Exploitation & 1,52 & 0,97 & 0,55 & 64~\% & Surtout ignorance \\
\end{tableaucours}
\vspace{-0.8em}
\begin{center}{\small\itshape Tableau~8.3. Décomposition de l'entropie de
trois domaines}\end{center}

\begin{boxlecture}{Deux entropies voisines, deux démarches}
D1 et D2 ont des entropies proches (1,41 et 1,30), pour des raisons
opposées.

\medskip
\textbf{Analyse} (confiance élevée). Dans D1, près des deux tiers de
l'entropie viennent de l'ignorance~: personne ne sait vraiment qui
désactive les comptes~; on cherche la preuve. Dans D2, la dispersion vient
du désaccord~: le DAF croit les RIB contrôlés, le chef de la liquidation
sait qu'ils ne le sont pas~; on confronte les deux versions à la trace. Le
même chiffre appelle deux démarches.

\medskip
\hfill\lecture{0{,}55}{0{,}30}{1{,}41}
\end{boxlecture}

\separateur

% ============================================================================
\section{Le quadrant conformité × incertitude}
\label{sec:c08-quadrant}
% ============================================================================

L'article A1 propose de lire ensemble la conformité et l'incertitude dans un
quadrant à quatre régions, chacune appelant une action
différente~\cite{minka2026nek}. Le cours retient, par convention, un seuil de
0,5 pour $\mu$ et de 0,20 pour $U$~; chaque mission peut fixer les siens, à
condition de les écrire dans le programme de travail.

\noindent\begin{minipage}{\linewidth}
\begin{center}
\begin{tikzpicture}[scale=0.95]
\fill[green!18] (3,0) rectangle (6,2);
\fill[red!15] (0,0) rectangle (3,2);
\fill[orange!25] (3,2) rectangle (6,5);
\fill[orange!15] (0,2) rectangle (3,5);
\draw[thick, color=CouleurPrimaire] (0,0) rectangle (6,5);
\draw[dashed, color=CouleurPrimaire] (3,0) -- (3,5);
\draw[dashed, color=CouleurPrimaire] (0,2) -- (6,2);
\node[align=center, font=\footnotesize\bfseries] at (1.5,0.45) {Non-conformité claire};
\node[align=center, font=\footnotesize\bfseries] at (4.5,0.45) {Conformité claire};
\node[align=center, font=\footnotesize\bfseries] at (1.5,4.6) {Non-conformité opaque};
\node[align=center, font=\footnotesize\bfseries] at (4.5,4.6) {Conformité opaque};
\foreach \x/\y/\t in {3.3/3.0/D1, 1.8/1.0/D2, 4.5/1.0/D3, 2.4/4.0/D4}
    \fill[CouleurAccent] (\x,\y) circle (2.5pt) node[above right, font=\scriptsize\bfseries] {\t};
\node[font=\scriptsize] at (3,-0.3) {0,5};
\node[font=\scriptsize] at (-0.35,2) {0,20};
\node[font=\small] at (3,-0.8) {Taux de conformité $\mu$};
\node[font=\small, rotate=90] at (-0.9,2.5) {Taux d'incertitude $U$};
\end{tikzpicture}
\end{center}
\vspace{-0.4em}
\begin{center}{\small\itshape Figure~8.1. Le quadrant diagnostique et les
quatre domaines du MINÉQUIR (d'après~\cite{minka2026nek}~; seuils du cours~;
échelle verticale $U \times 10$)}\end{center}
\end{minipage}

\medskip

\begin{tableaucours}{L{3.4cm}L{4.0cm}L{4.8cm}}{%
    \tableauHeader{Région} &
    \tableauHeader{Situation} &
    \tableauHeader{Action}
}
Conformité claire & On s'accorde, avec assurance, sur la conformité &
Validation courante \\
Non-conformité claire & On s'accorde, avec assurance, sur la
non-conformité & Correctifs \\
\textbf{Conformité opaque} & Des \OUI{} nombreux, mais une part importante
de \NEK{} & \textbf{Alerte prioritaire}~: rechercher la preuve avant toute
conclusion \\
Non-conformité opaque & Des \NON{} nombreux et une part importante de
\NEK{} & Rechercher la preuve et engager des premiers correctifs \\
\end{tableaucours}
\vspace{-0.8em}
\begin{center}{\small\itshape Tableau~8.4. Les quatre régions et leurs
actions (d'après~\cite{minka2026nek})}\end{center}

La conformité opaque est la région la plus dangereuse~: une grille binaire
la confond avec la conformité claire. C'est le domaine D1 du MINÉQUIR, celui
du double paiement. L'article A0 dit la même chose dans son
vocabulaire~: la \emph{non-conformité connue} appelle des correctifs~; la
\emph{non-conformité inconnue}, que cachent les régions opaques, appelle
d'abord des investigations~\cite{minka2026auditologie}.

\separateur

% ============================================================================
\section{La règle de décision du cours}
\label{sec:c08-regle}
% ============================================================================

\begin{boxprocedure}
\textbf{Conclure sur un domaine.}
\begin{enumerate}
    \item \textbf{$U$ d'abord.} Si $U \geq 0{,}20$, le domaine est opaque~:
          aucune conclusion de conformité n'est possible. L'auditeur
          recherche la preuve, ou conclut \emph{sur l'opacité elle-même}.
    \item \textbf{$\mu$ ensuite.} Si $U < 0{,}20$, le domaine est clair~:
          conforme si $\mu \geq 0{,}5$, non conforme sinon.
    \item \textbf{$H$ pour contrôler.} Une entropie élevée avec $U$ faible
          signale un désaccord~: le domaine n'est pas homogène~; on le
          décompose avant de conclure (Chapitre~\ref{chap:preparer}).
    \item \textbf{Les cotes de l'auditeur prévalent.} Les mesures des
          répondants orientent~; la conclusion repose sur les contrôles
          cotés par l'auditeur, preuve à l'appui (Chapitre~\ref{chap:preuve}).
\end{enumerate}
\end{boxprocedure}

\begin{boxlecture}{Les quatre domaines du MINÉQUIR}
\textbf{Analyse} (confiance moyenne).
\begin{itemize}
    \item \emph{D1}~: conformité opaque~; alerte prioritaire. Les tests
          (Chapitre~\ref{chap:acces}) établiront un seul contrôle sur six
          conforme~: l'opacité cachait une non-conformité.
    \item \emph{D2}~: non-conformité claire~; correctifs immédiats sur les
          RIB.
    \item \emph{D3}~: conformité claire selon les répondants~; mais
          l'auditeur a coté \NON{} le blocage des doublons (règle~4)~: la
          conclusion suit les cotes.
    \item \emph{D4}~: non-conformité opaque~; investigations et premiers
          correctifs.
\end{itemize}
\end{boxlecture}

\begin{boxalert}
\textbf{Violation typique.} Moyenner les quatre domaines pour donner une
note globale («~le circuit est conforme à 50~\%~»).

\smallskip
\textbf{Règle~:} une moyenne efface les régions du quadrant. Le rapport
présente chaque domaine dans sa région.
\end{boxalert}

\separateur

% ============================================================================
\section{Le constat en cinq éléments}
\label{sec:c08-constat}
% ============================================================================

\subsection{Le constat}

\begin{tableaucours}{L{2.6cm}L{4.4cm}L{5.2cm}}{%
    \tableauHeader{Élément} &
    \tableauHeader{Question} &
    \tableauHeader{Règle de rédaction}
}
Condition & Qu'a-t-on constaté~? & Un fait vérifiable, chiffré, daté, avec
sa source et sa cote \\
Critère & Que fallait-il~? & La règle citée exactement, avec sa force
(Chapitre~\ref{chap:normes}) \\
Cause & Pourquoi l'écart~? & Ce qui, dans le dispositif, l'explique~; la
cause appelle la recommandation \\
Effet & Qu'est-ce que cela produit ou permet~? & Le risque réalisé ou
possible, chiffré quand on le peut \\
Recommandation & Que faire~? & Une action, un responsable, une échéance \\
\end{tableaucours}
\vspace{-0.8em}
\begin{center}{\small\itshape Tableau~8.5. Les cinq éléments d'un
constat}\end{center}

\subsection{L'observation en sept éléments}

La méthode du cours d'audit de l'auteur développe le constat en une
\termecle{observation}\index{observation} en sept éléments. Les deux
éléments ajoutés sont indispensables à l'opposabilité (axiome~6
d'A0~\cite{minka2026auditologie})~: la trace du contradictoire, et une
conclusion explicite qui sépare la position de l'auditeur de son analyse.

\begin{tableaucours}{L{5.0cm}L{6.0cm}}{%
    \tableauHeader{Constat en cinq éléments} &
    \tableauHeader{Observation en sept éléments}
}
Critère & 1. Norme \\
Condition & 2. Constatation (avec la cote) \\
Effet & 3. Risques \\
(absent) & 4. Contradictoire \\
Cause & 5. Analyse (avec un niveau de confiance) \\
(absent) & 6. Conclusion \\
Recommandation & 7. Recommandations \\
\end{tableaucours}
\vspace{-0.8em}
\begin{center}{\small\itshape Tableau~8.6. Du constat à
l'observation}\end{center}

\begin{boxobjet}{D1.1.2 · Observation}
\textbf{1. Norme~:} le compte d'un agent est désactivé à la date d'effet de
sa mutation (bonne pratique~: COSO~2013, principe~11~; aucune règle
interne).

\textbf{2. Constatation~:} au 23~mars~2026, 9 des 31 agents mutés ou partis
en 2025 ont un compte actif, dont \texttt{mbarga.j}, qui cumule LIQ, ORD et
FRS (extractions et empreintes au dossier). Cote \NON{}.

\textbf{3. Risques~:} un compte resté actif a servi, les 27 et
28~décembre~2025, à modifier un RIB et à émettre seul un mandat de
187\,450\,000~FCFA.

\textbf{4. Contradictoire~:} le 30~avril~2026, la DAF indique que les
mutations ne lui sont pas notifiées par les ressources humaines.

\textbf{5. Analyse} (confiance élevée)~: la cause est l'absence de flux
entre le fichier du personnel et la cellule informatique, et l'absence de
procédure~; non la négligence d'un agent.

\textbf{6. Conclusion~:} le contrôle D1.1.2 n'est pas efficace~;
déficience majeure du domaine D1 (Chapitre~\ref{chap:eval-sci}).

\textbf{7. Recommandations~:} désactiver les 9 comptes~; organiser la
notification écrite des mutations par la DRH~; imposer dans SIDÉPENSE
l'incompatibilité LIQ~/ ORD.
\end{boxobjet}

\begin{boxalert}
\textbf{Violation typique.} Écrire dans les risques~: «~M.~MBARGA a détourné
187\,450\,000~FCFA~».

\smallskip
\textbf{Règle~:} l'observation dit ce que la preuve établit~: un compte, pas
une personne. L'imputation relève d'une autre procédure.
\end{boxalert}

\subsection{Écrire une recommandation}

\begin{tableaucours}{L{5.6cm}L{6.6cm}}{%
    \tableauHeader{Mal écrite} &
    \tableauHeader{Bien écrite}
}
Renforcer la sécurité des accès & Organiser la notification écrite des
mutations de la DRH vers la cellule, avant le 31~juillet~2026 (responsable~:
DRH) \\
Désactiver les comptes inutiles & Désactiver sous 48~heures les 9 comptes
identifiés (annexe~C), puis revoir chaque semestre les comptes actifs
(responsable~: DAF) \\
Sanctionner M.~MBARGA & (Hors du rôle de l'auditeur~: la sanction relève
de l'autorité compétente, après enquête.) \\
\end{tableaucours}
\vspace{-0.8em}
\begin{center}{\small\itshape Tableau~8.7. Écrire une
recommandation}\end{center}

\separateur

% ============================================================================
\section{Le rapport}
\label{sec:c08-rapport}
% ============================================================================

\begin{enumerate}
    \item \textbf{Synthèse}~: l'objectif, la conclusion sur l'efficacité du
          contrôle interne, les domaines dans leur région du quadrant, les
          déficiences majeures, en une page.
    \item \textbf{Mandat et étendue}~: fondement, périmètre, période,
          travaux, seuils de décision, et \termecle{limitations
          d'étendue}\index{limitation d'étendue} (les \NEK{} restants, avec
          les diligences faites).
    \item \textbf{Mesures par domaine}~: $\mu$, $U$, $H$ et région.
    \item \textbf{Observations}, classées par domaine puis par gravité.
    \item \textbf{Tableau des déficiences} et leur classement.
    \item \textbf{Réponses de l'entité} et position de l'auditeur.
    \item \textbf{Annexes}~: matrice des risques et des contrôles,
          décomposition, liste des pièces et de leurs empreintes.
\end{enumerate}

Un \NEK{} qui reste n'est pas caché~: la restitution impartiale l'exige, et
la conclusion ne peut pas porter sur ce qui n'a pas pu être établi.

\separateur

% ============================================================================
\section{La procédure contradictoire}
\label{sec:c08-contradictoire}
% ============================================================================

\subsection{Pourquoi le contradictoire}

Le rapport provisoire est adressé à l'entité, qui répond par écrit dans un
délai fixé. Le contradictoire n'est pas une politesse~: il réduit
l'incertitude. L'audité connaît des faits que l'auditeur ignore~; sa
réponse peut lever un \NEK{}, corriger une cause, ou révéler une pièce. Il
est aussi la condition de l'opposabilité~: une conclusion que l'audité n'a
pas pu discuter se défend mal devant le CDBF ou la Chambre des Comptes.

\begin{tableaucours}{L{4.4cm}L{4.0cm}L{3.8cm}}{%
    \tableauHeader{Réponse de la DAF (30 avril 2026)} &
    \tableauHeader{Position de la vérificatrice} &
    \tableauHeader{Effet sur le rapport}
}
«~Les mutations ne nous ont pas été notifiées par les ressources
humaines.~» & Exact~; c'est la cause relevée~; ne change pas la condition
& Observation maintenue, cause précisée \\
«~Les 9 comptes ont été désactivés le 20~avril.~» & Pièce fournie
(journal) & Mention au rapport~; R-01 mise en œuvre \\
«~La séparation a toujours été respectée en pratique.~» & Contredit par
le mandat 2025-09342 (journal) & Observation maintenue \\
\end{tableaucours}
\vspace{-0.8em}
\begin{center}{\small\itshape Tableau~8.8. Le contradictoire du dossier
MINÉQUIR}\end{center}

\subsection{Maintenir, modifier, retirer}

Pour chaque réponse, l'auditeur décide, et motive~: il \emph{maintient}
l'observation, la \emph{modifie} au vu d'une pièce nouvelle, ou la
\emph{retire}. Une déclaration sans pièce ne fait jamais changer une cote
\NON{}~; elle peut, au mieux, transformer une cause.

\begin{boxprocedure}
\textbf{Fiche de contradictoire}, une par observation~: date d'envoi du
rapport provisoire~; réponse reçue (date, signataire, texte)~; pièces
nouvelles~; position motivée~; modification apportée au rapport définitif.
\end{boxprocedure}

\separateur

% ============================================================================
\section{Suivre les recommandations}
\label{sec:c08-suivi}
% ============================================================================

\begin{tableaucours}{cL{5.0cm}L{2.6cm}L{2.0cm}L{1.8cm}}{%
    \tableauHeader{N°} &
    \tableauHeader{Recommandation} &
    \tableauHeader{Responsable} &
    \tableauHeader{Échéance} &
    \tableauHeader{Statut}
}
R-01 & Désactiver les 9 comptes & Cellule informatique & 20/04/2026 &
Fait \\
R-02 & Imposer l'incompatibilité LIQ / ORD dans SIDÉPENSE & Cellule,
éditeur & 30/06/2026 & En cours \\
R-03 & Double validation des RIB & DAF, cellule & 30/06/2026 & En cours
\\
R-04 & Bloquer les doublons (fournisseur, montant, facture normalisée) &
Cellule, éditeur & 30/06/2026 & En cours \\
R-05 & Notification des mutations par les ressources humaines & DRH &
31/07/2026 & Non commencé \\
R-06 & Revue des mandats de fin de gestion & DAF, contrôle financier &
31/12/2026 & Non commencé \\
\end{tableaucours}
\vspace{-0.8em}
\begin{center}{\small\itshape Tableau~8.9. Plan d'action du MINÉQUIR
(situation au 15~mai~2026)}\end{center}

Le suivi vérifie la mise en œuvre par une preuve, pas par une
déclaration~: pour R-02, la vérificatrice demandera une tentative, en
environnement de test, d'attribuer LIQ et ORD au même compte.

\begin{boxinfo}
\textbf{L'audit modifie ce qu'il observe.} L'article A0 énonce une
proposition de \emph{réflexivité}~: l'acte d'audit fait partie de la
trajectoire du système audité~\cite{minka2026auditologie}. Le plan d'action
change le système~; le cycle suivant n'audite donc pas le même objet. C'est
pourquoi on compare des distributions, et non des taux isolés.
\end{boxinfo}

\separateur

% ============================================================================
\section{Comparer deux cycles}
\label{sec:c08-cycles}
% ============================================================================

\subsection{Pourquoi le taux de conformité ne suffit pas}

Le taux de conformité seul peut rester stable alors que l'état de
connaissance change~: des \NEK{} deviennent des \NON{}, l'entité sait
désormais ce qu'elle ignorait. Il faut comparer les \emph{distributions} de
réponses. C'est la conséquence de l'\emph{individuation temporelle}
(axiome~5)~: les conclusions d'une mission ne se testent que sur plusieurs
cycles~\cite{minka2026auditologie}.

\begin{boxdef}
\termecle{Divergence de Kullback-Leibler}\index{divergence de
Kullback-Leibler}~: pour deux distributions $p$ (nouveau cycle) et $q$
(cycle précédent) sur $\{\OUI, \NON, \NEK\}$~\cite{kullback1951,minka2026nek},
\[
    D_{\mathrm{KL}}(p \,\|\, q) = \sum_{m} p(m) \log_2 \frac{p(m)}{q(m)}.
\]
Elle est positive ou nulle, nulle si et seulement si les deux distributions
sont identiques~; elle n'est pas symétrique (on mesure toujours le nouveau
cycle par rapport à l'ancien)~; elle s'exprime en bits.
\end{boxdef}

\begin{boxscenario}{Dossier MINÉQUIR. Le domaine D1 en novembre 2026}
\textbf{Faits.} En mars~2026, D1 donnait 11~\OUI{}, 3~\NON{}, 6~\NEK{}. À
l'audit de suivi de novembre~2026, après R-01 à R-03, les mêmes questions
posées aux mêmes fonctions donnent 17~\OUI{}, 2~\NON{}, 1~\NEK{}.

\medskip
\textbf{Calcul.} $\mu$ passe de 0,55 à 0,85, $U$ de 0,30 à 0,05, $H$ de 1,41
à 0,75~bit. La divergence vaut
\[
    0{,}85\log_2\frac{0{,}85}{0{,}55} + 0{,}10\log_2\frac{0{,}10}{0{,}15} +
    0{,}05\log_2\frac{0{,}05}{0{,}30} \approx 0{,}534 - 0{,}058 - 0{,}129
    \approx 0{,}35~\text{bit}.
\]

\textbf{Analyse} (confiance moyenne). D1 est passé de la conformité opaque
à la conformité claire. La vérificatrice confirmera par ses propres tests
(revue des comptes actifs, essai de cumul en environnement de test) avant
de conclure.
\end{boxscenario}

\subsection{Le lissage}

Si une réponse est absente au premier cycle et présente au second, un terme
de la divergence devient infini~: ce n'est pas que le changement soit
infini, c'est que le premier échantillon, petit, n'a rien observé. On
ajoute alors une unité fictive à chaque réponse des deux cycles (lissage de
Laplace). La divergence de D1 devient 0,25~bit. On applique le même lissage
à toutes les comparaisons d'une mission, et on l'écrit dans le programme.

\begin{boxoutils}{Calculer la divergence (Python)}
\begin{lstlisting}[language=Python]
from math import log2
def dkl(nouveau, ancien, lissage=0):
    a = [x + lissage for x in nouveau]
    b = [x + lissage for x in ancien]
    p = [x / sum(a) for x in a]
    q = [x / sum(b) for x in b]
    return sum(pi * log2(pi / qi) for pi, qi in zip(p, q) if pi > 0)

print(round(dkl([17, 2, 1], [11, 3, 6]), 3))     # 0.346
print(round(dkl([17, 2, 1], [11, 3, 6], 1), 3))  # 0.247 (lissée)
\end{lstlisting}
\end{boxoutils}

\separateur

% ============================================================================
\section{Quand le constat devient une affaire}
\label{sec:c08-suites}
% ============================================================================

Une observation peut fonder une suite~: saisine du CDBF, que le CONSUPE peut
engager (décret n°~2013/287, art.~3 et~5)~; transmission à la Chambre des
Comptes~; plainte pénale. L'auditeur n'instruit pas ces procédures, mais il
les rend possibles ou impossibles par la qualité de sa preuve.

\begin{itemize}
    \item \textbf{Geler la preuve}~: conserver les extractions originales,
          leurs empreintes et les procès-verbaux~; travailler sur des copies.
    \item \textbf{Ne pas alerter}~: ne pas interroger la personne soupçonnée
          sur les faits avant que l'autorité ait décidé de la suite.
    \item \textbf{Demander la conservation} des journaux au-delà de leur
          durée habituelle.
    \item \textbf{Séparer}~: le rapport traite du contrôle interne~; une note
          distincte, adressée à la hiérarchie, présente les faits
          susceptibles de poursuite.
\end{itemize}

\begin{boxmetier}{Vérificateur}
Présentez dans la synthèse les domaines par région du quadrant, et non par
une note globale. Une «~conformité opaque~» écrite en toutes lettres
protège votre rapport mieux qu'un taux flatteur.
\end{boxmetier}

\begin{boxmetier}{Juge des comptes}
Une limitation d'étendue écrite au rapport dit ce que l'auditeur n'a pas pu
établir~; elle désigne les pièces que l'instruction peut exiger.
\end{boxmetier}

\begin{boxmetier}{Auditeur interne}
Refaites chaque année les mêmes questions aux mêmes fonctions, et calculez
la divergence entre les cycles~: c'est la mesure la plus honnête de
l'effet de vos recommandations.
\end{boxmetier}

\separateur

\begin{boxresume}
\begin{itemize}
    \item \textbf{Trois mesures}~: $\mu$ (conformité), $U$ (incertitude), $H$
          (entropie, au plus $\log_2 3$)~; jamais l'une sans les autres.
    \item \textbf{L'entropie se décompose} en ignorance et désaccord~; chacun
          appelle sa démarche.
    \item \textbf{Le quadrant} sépare quatre régions~; la conformité opaque
          est la plus dangereuse.
    \item \textbf{La règle de décision}~: $U$ d'abord, $\mu$ ensuite, $H$
          pour contrôler~; les cotes de l'auditeur prévalent.
    \item \textbf{Sept éléments}~: l'observation ajoute au constat le
          contradictoire et la conclusion, pour l'opposabilité.
    \item \textbf{Le contradictoire réduit l'incertitude}~; une déclaration
          sans pièce ne fait pas changer un \NON{}.
    \item \textbf{Comparer les cycles} par la divergence de Kullback-Leibler,
          avec le même lissage partout.
    \item \textbf{Préserver la suite}~: geler la preuve, ne pas alerter,
          séparer rapport et note de faits.
\end{itemize}
\end{boxresume}

\begin{boxexo}[title={\ding{45}\quad Exercice~8.1. Calculer et situer \niveaubase}]
Calculez $\mu$, $U$ et $H$ pour (a) 18~\OUI{}, 1~\NON{}, 1~\NEK{}~; (b)
0~\OUI{}, 0~\NON{}, 20~\NEK{}~; (c) 10~\OUI{}, 8~\NON{}, 2~\NEK{}. Situez
chacun dans le quadrant et dites l'action.
\end{boxexo}

\begin{boxexo}[title={\ding{45}\quad Exercice~8.2. Ignorance ou désaccord~? \niveaumoyen}]
Décomposez l'entropie du domaine D3 (15~\OUI{}, 3~\NON{}, 2~\NEK{}) et
dites quelle démarche elle appelle. Expliquez pourquoi la conclusion du
domaine reste \NON{} sur le blocage des doublons.
\end{boxexo}

\begin{boxexo}[title={\ding{45}\quad Exercice~8.3. L'observation \niveaumoyen}]
Rédigez l'observation en sept éléments sur le contrôle «~seconde validation
des modifications de RIB~» (domaine D2), à partir des faits du prologue et
de la réponse de la DAF~: «~les RIB sont modifiés sur demande écrite du
fournisseur~» (aucune demande n'est produite pour BATIRURAL).
\end{boxexo}

\begin{boxexo}[title={\ding{45}\quad Exercice~8.4. Deux cycles \niveauexpert}]
À l'audit de suivi, D4 donne 8~\OUI{}, 10~\NON{}, 2~\NEK{} (contre 8, 4, 8
en mars). (a) Le taux de conformité a-t-il changé~? (b) Calculez la
divergence. (c) Dites si l'exploitation s'est améliorée, et rédigez la
phrase du rapport de suivi.
\end{boxexo}

\begin{boxinfo}
\textbf{Transition.} La méthode est complète, du mandat au cycle suivant. La
Partie~IV l'applique domaine par domaine, en commençant par celui dont tout
dépend~: les accès (Chapitre~\ref{chap:acces}).
\end{boxinfo}
